\runninghead{Each Text in Its Setting}
\section{Each Text in Its Setting}
\label{sec:each-element}

\subsection{Introit: a letter to the exiles}
\label{sec:element-introit}

The antiphon's words come from a letter that Jeremias sent from Jerusalem
``to the residue of the ancients that were carried into captivity, and to the
priests, and to the prophets, and to all the people, whom Nabuchodonosor had
carried away from Jerusalem to Babylon'' (Jer 29:1). The letter does not tell
the exiles to expect a quick return. They are to ``Build ye houses, and dwell
in them: and plant orchards, and eat the fruit of them'' (v.~5), to ``seek the
peace of the city, to which I have caused you to be carried away captives''
(v.~7), and to give no heed to the prophets and diviners among them, who
``prophesy falsely to you in my name'' (vv.~8--9). Only after that does God
promise to visit them ``When the seventy years shall begin to be accomplished
in Babylon'' (v.~10), and the words that the Introit sings follow
(vv.~11--14).

The antiphon keeps three things that God says he will do: think thoughts of
peace, hear those who call on him, and bring back their captivity. Its
opening, \latin{Dicit Dóminus}, makes the whole of it God's speech in the first
person, as it is in the letter. St Jerome comments on the first of these words
by setting God's knowledge against the exiles' prophets: \latin{Se dicit nosse
quid cogitet, illos autem cum suis prophetis \dots\ ignorare. Futurorum igitur
scientia soli Deo competit}, ``He says that he knows what he thinks, and that
they, with their prophets, do not know. The knowledge of things to come
therefore belongs to God alone'' (\work{In Hieremiam}~V, on 29:10--14, PL~24
col.~893). The antiphon leaves out the promise ``to give you an end and
patience,'' the seeking ``with all your heart'' (v.~13) and the gathering
``out of all nations.'' Its last words, ``from all places,'' stand where the
letter promises to gather the exiles ``from all the places to which I have
driven you out'' (v.~14).

The verse, Ps 84:2, opens a psalm ``for the sons of Core'' (84:1) whose first
verses speak in past tenses: ``Lord, thou hast blessed thy land: thou hast
turned away the captivity of Jacob. Thou hast forgiven the iniquity of thy
people: thou hast covered all their sins'' (vv.~2--3). Two verses later the
same psalm prays, ``Convert us, O God our saviour: and turn off thy anger from
us'' (v.~5). The antiphon promises a return in the future tense, and the verse
sung with it speaks of the captivity as already turned away.

\subsection{Collect: wills stirred to the fruit of the divine work}
\label{sec:element-collect}

The Collect opens with an imperative, \latin{Excita}, ``stir up'' or
``arouse,'' and its object is \latin{tuórum fidélium voluntátes}, ``the wills
of Thy faithful'' in Lasance's English. The rest of the prayer gives the
purpose: that, ``by more earnestly following after the fruit of the divine
work, they may the more abundantly partake of Thy mercies.'' The Latin sets
two comparatives side by side. The faithful follow the fruit
\latin{propénsius}, ``more readily,'' and receive \latin{remédia}, ``remedies,''
that are \latin{maióra}, ``greater''; Lasance's ``mercies'' stands for those
remedies. \latin{Divíni óperis}, ``of the divine work,'' can be the work that
God does or the work that is done for God, and the prayer does not say
which. The Epistle that follows prays that the Colossians may be
filled with the knowledge of ``the will of God'' and be ``fruitful in every
good work''; the Collect uses the same two words, will and fruit.

\subsection{Epistle: Paul's prayer for a church he had not founded}
\label{sec:element-epistle}

Paul writes to the Colossians, with Timothy, after hearing of their ``faith in
Christ Jesus and the love which you have towards all the saints'' (Col 1:4)
from Epaphras, ``our most beloved fellow servant,'' who had taught them
(vv.~7--8). He has already given thanks for them (v.~3). The lesson begins in
the middle of the next sentence, where thanks turns to petition, and the
Missal sets ``Brethren'' in place of Paul's ``Therefore we also, from the day
that we heard it.''

The lesson is a single sentence. Its petition is that they ``may be filled with
the knowledge of the will of God, in all wisdom and spiritual understanding''
(v.~9), and its purpose that they may ``walk worthy of God, in all things
pleasing'' (v.~10). Four participles describe that walk: bearing fruit in
every good work, growing in the knowledge of God, strengthened with all might
``in all patience and longsuffering with joy,'' and giving thanks (vv.~10--12).
The thanks are given to the Father, and three relative clauses say what he has
done: he ``hath made us worthy to be partakers of the lot of the saints in
light,'' he ``hath delivered us from the power of darkness,'' and he ``hath
translated us into the kingdom of the Son of his love'' (vv.~12--13). The last
clause names the Son: ``In whom we have redemption through his blood, the
remission of sins'' (v.~14). The petition is for something not yet complete,
and the clauses of thanks speak of what has already been done.

St John Chrysostom notices both. On ``that ye may be filled'' he writes that
Paul ``nowhere says that they are destitute of everything, but that they are
deficient,'' and that he says ``\,`that ye may be filled'; not, `that ye may
receive,' for they had received'' (\work{Hom.\ in Col.}~2). On ``the lot of
the saints'' he asks why Paul calls it a lot, and answers ``To show that by
his own achievements no one obtains the kingdom,'' since ``the whole is of His
free gift'' (Lk 17:10). On the light he says that Paul ``means, both the future
and the present light.''

\subsection{Gradual: thanks from a psalm of complaint}
\label{sec:element-gradual}

Psalm 43 begins with the remembrance of what God did for the fathers: ``For
they got not the possession of the land by their own sword: neither did their
own arm save them. But thy right hand and thy arm'' (v.~4). It goes on in the
first person, ``For I will not trust in my bow,'' and then gives the two verses
that the Gradual sings: ``But thou hast saved us from them that afflict us:
and hast put them to shame that hate us. In God shall we glory all the day
long'' (vv.~7--9). The next verse turns: ``But now thou hast cast us off, and
put us to shame: and thou, O God, wilt not go out with our armies'' (v.~10).
The rest of the psalm is a complaint of a people given up ``like sheep to be
eaten'' (v.~12), though ``we have not forgotten thee'' (v.~18), and near its
end it cries, ``Arise, why sleepest thou, O Lord?'' (v.~23). The Gradual sings
the confidence and not the complaint.

St Augustine, preaching on the psalm, takes the complaint for the voice of
Christians suffering persecution, and asks on their behalf, \latin{Quare nos
ista omnia medio tempore pati voluit?}, ``Why did he will us to suffer all
these things in the meantime?'' (\work{En.\ in Ps.}~43.10). What he says of
the Gradual's past tense is given in the third interpretation.

\subsection{Alleluia: from the depths}
\label{sec:element-alleluia}

Psalm 129 is one of the gradual canticles. Its eight verses go from the cry
to hope: ``If thou, O Lord, wilt mark iniquities: Lord, who shall stand it.
For with thee there is merciful forgiveness \dots\ Because with the Lord there
is mercy: and with him plentiful redemption. And he shall redeem Israel from
all his iniquities'' (vv.~3--4, 7--8). The Alleluia sings only the cry and the
request to be heard, in the Missal's form ``my prayer'' (vv.~1--2), and the
same words return at the Offertory.

St Hilary of Poitiers asks the reader to picture the psalmist \latin{de
profundo humanarum calamitatum aestu precantem}, ``praying from the deep surge
of human calamities,'' and to suppose him praying \latin{non labiis, sed
secreto \dots\ cordis}, ``not with the lips, but in the secret of the heart,''
as the Lord commanded (Mt 6:6). The cry, he says, fits ignorance, suffering or
prayer, \latin{dum aut nescit, aut patitur, aut precatur}, ``while he either
does not know, or suffers, or prays'' (\work{Tractatus super Psalmos}, on Ps
129, PL~9 col.~718).

\subsection{Gospel: the second part of the discourse on the Mount of Olivet}
\label{sec:element-gospel}

The disciples asked two things: when the temple would be destroyed, and what
would be the sign of the Lord's coming ``and of the consummation of the
world'' (Mt 24:3). The first part of the answer, which the Mass does not read,
warns against those who will come saying ``I am Christ,'' speaks of wars,
famines and persecution as ``the beginnings of sorrows,'' and ends ``this
gospel of the kingdom shall be preached in the whole world \dots\ and then
shall the consummation come'' (vv.~4--14). The Missal's reading begins with
the next verse and drops its ``therefore.''

The passage falls into four parts. Verses 15--22 give the sign from Daniel
(Dan 9:27) and the commands to flee, and close with the days shortened for the
elect. Verses 23--28 warn against false christs who will be said to be ``in
the desert'' or ``in the closets,'' set the lightning against them, and end
with the body and the eagles. Verses 29--31 describe the darkened sky, the
sign of the Son of man, his coming on the clouds and the gathering of the
elect. Verses 32--35 give the parable of the fig tree, the saying about this
generation and the saying about heaven and earth. The discourse continues with
``But of that day and hour no one knoweth'' (v.~36), which the Mass does not
read.

The parenthesis ``he that readeth let him understand'' is the first
difficulty. St Jerome takes it as a signal: \latin{Quando ad intelligentiam
provocamur, mysticum monstratur esse quod dictum est}, ``When we are summoned
to understanding, what has been said is shown to be mystical'' (\work{In
Matthaeum}~IV, on 24:15). He then gives three plain senses of the
abomination, \latin{aut de Antichristo \dots\ aut de
imagine Caesaris, quam Pilatus posuit in templo; aut de Hadriani equestri
statua}, ``either of Antichrist, or of the image of Caesar which Pilate placed
in the temple, or of the equestrian statue of Hadrian,'' and does not choose
among them. St John Chrysostom takes it to be ``the statue of him who then
took the city,'' set up in the temple, and says the Lord spoke of it ``in order
that they might learn that these things will be while some of them are alive''
(\work{Hom.\ in Mt.}~75.2). St Thomas Aquinas mentions the
Roman eagle that Pilate brought into the temple and the idol of Hadrian
(\work{Super Matthaeum}, cap.~24).

The second difficulty is the order of time. The word ``Then'' at verse~23
seems to follow directly on the flight from Judea. Chrysostom says that it
``relates not to the connection in order of time,'' and that the Lord here
passes over ``all the intermediate time from the taking of Jerusalem unto the
preludes of the consummation'' (\work{Hom.\ in Mt.}~76.2). St Jerome, on the
false christs of the same verse, allows that they may be those who rose during
the Roman siege, \latin{Sed melius de consummatione mundi intelligitur}, ``but
it is better understood of the consummation of the world'' (on 24:23). Both
read verses 29--31 of the Lord's visible coming. The three interpretations
below take up, in turn, the eagles of verse 28 and the gathering of verse 31,
the flight of verses 15--20, and the fig tree and the two sayings of verses
32--35.

\subsection{Offertory: the cry sung again}
\label{sec:element-offertory}

The Offertory sings the Alleluia's words again while the gifts are offered,
and returns at the end to its first line, so that the cry ``From the depths''
is sung three times in the Mass. Rupert of Deutz and Sicard of Cremona,
medieval writers on the liturgy, expound the Introit, Gradual, Offertory and
Communion in turn; what they say of the Gradual, the Offertory and the
Communion is given in the first and third interpretations.

\subsection{Secret: hearts turned and desires changed}
\label{sec:element-secret}

The Secret asks God, in Lasance's English, to be ``propitious \dots\ to our
supplications,'' to ``accept the offerings and prayers of Thy people,'' and to
``turn all our hearts unto Thee, that, being delivered from earthly desires,
we may pass on to the enjoyments of heaven.'' The Latin's purpose clause has a
participle, \latin{liberáti}, ``delivered,'' and a verb of motion,
\latin{transeámus}, ``may we pass over.'' Its \latin{liberáti} is the verb the Gradual uses of what God has done,
\latin{Liberásti nos}, ``Thou hast delivered us.''

\subsection{Communion: the promise after the withered fig tree}
\label{sec:element-communion}

The setting of the Communion's saying is given in the opening: the fig tree
cursed and withered, Peter's remark, and the Lord's answer, ``Have the faith of
God'' (Mk 11:20--22). The verse before the Communion's promises that whoever
says to this mountain ``Be thou removed and be cast into the sea, and shall not
stagger in his heart, but believe that whatsoever he saith shall be done; it
shall be done unto him'' (v.~23). That verse begins ``Amen I say to you'' and
ends with ``it shall be done,'' the words with which the antiphon opens and
closes, where it departs from the Vulgate's verse~24. The promise is in the future tense, \latin{crédite quia
accipiétis}, ``believe that you shall receive,'' and the next verse adds a
condition to it: ``when you shall stand to pray, forgive, if you have aught
against any man'' (v.~25).

St Bede, on the mountain-moving faith of verse 23, answers pagans who say
that Christians have never moved a mountain. Not everything done in the
Church is written, he says, and a mountain could have been cast into the sea
if need had required it. He tells how a mountain gave way at the prayers of
Gregory, bishop of Neocaesarea, as far as was needed to build a church. He
also takes the mountain for the devil, driven out of hearts when holy teachers
preach the word (\work{In Marci evangelium expositio}~III, on
11:22--25).

\subsection{Postcommunion: healing of the mind}
\label{sec:element-postcommunion}

The Postcommunion asks that, ``through this sacrament which we have received,
whatever is evil in our hearts may be restored by its gift of healing''
(Lasance). Its Latin names the cure \latin{medicatiónis}, ``of healing,'' and
the Collect asked for \latin{remédia}, ``remedies,'' so both prayers speak of
cure. Dom Lucien Fromage, the Solesmes monk who continued Dom
Guéranger's \work{The Liturgical Year}, prays that the Sacrament may cure
``whatsoever there may remain faulty in our souls, at this close of the Year''
(vol.~11, p.~513).
